\content{
        \begin{example} Your turn! Use the Plurality with Elimination method to
        determine who wins. \\
        \begin{tabular}{|l|c|c|c|c|c|}
        \hline
        & 44 & 14 & 70 & 22 & 80 \\
        \hline
        1st & G&G&M&M&B\\
        \hline
        2nd & M&B&G&B&M\\
        \hline
        3rd & B&M&B&G&G\\
        \hline
        \end{tabular}
        \end{example}
}{% presentation
\vspace{3in}
}{% nopresentation
}{% comments
}

\content{
\begin{example} Use the plurality method with elimination to determine who wins
the election.\\
\begin{tabular}{|l|c|c|c|}
\hline
& 342 & 214 & 298 \\
\hline
1st & E&D&K\\
\hline
2nd & D&K&D\\
\hline
3rd & K&E&E\\
\hline
\end{tabular}

Is there a Condorcet Candidate? If so, did they win the election? 
\end{example}
}{% presentation
\vspace{3in}
}{% nopresentation
}{% comments
}

\content{
\begin{example} Who wins the following election under the plurality with
elimination method? \\
\begin{tabular}{|l|c|c|c|c|}
\hline
& 37 & 22 & 12 & 29 \\
\hline
1st & A&B&B&C\\
\hline
2nd & B&C&A&A\\
\hline
3rd & C&A&C&B\\
\hline
\end{tabular}
\end{example}
}{% presentation
\vspace{2in}
}{% nopresentation
}{% comments
}

\content{
Now, Look back at example 1.2.3, and suppose that after officials have announced the results, the ballots were
accidentally destroyed before they could be certified, so the votes have to be
recast. Ten of the people whose preference was $B,C,A$ decide that they want to
now support the winner $A$, so they change their votes to $A,B,C$, all the other
votes remain the same. What are the results of the election now? Start by
filling out what the preference schedule would look like after the new votes are
collected. \\
\begin{tabular}{|l|p{0.25in}|p{0.25in}|p{0.25in}|p{0.25in}|}
\hline
&  &  &  &  \\
\hline
1st & A&B&B&C\\
\hline
2nd & B&C&A&A\\
\hline
3rd & C&A&C&B\\
\hline
\end{tabular}
}{% presentation
\vspace{3in}\\

What happened? Is this fair? 
}{% nopresentation
}{% comments
}
